% Generated draft: requires mathematical and source review.
\documentclass[11pt]{article}
\usepackage[margin=25mm]{geometry}
\usepackage{fontspec}
\setmainfont{Latin Modern Roman}
\usepackage{amsmath,amssymb,mathtools,unicode-math}
\setmathfont{Latin Modern Math}
\usepackage{enumitem,needspace,xurl,hyperref}
\setlength{\parindent}{0pt}
\setlength{\parskip}{5pt}
\setlength{\emergencystretch}{3em}
\newcommand{\N}{\mathbb{N}}
\newcommand{\Z}{\mathbb{Z}}
\newcommand{\Q}{\mathbb{Q}}
\newcommand{\R}{\mathbb{R}}
\newcommand{\C}{\mathbb{C}}
\newcommand{\F}{\mathbb{F}}
\newcommand{\A}{\mathbb{A}}
\newcommand{\PP}{\mathbb{P}}
\newcommand{\E}{\mathbb{E}}
\newcommand{\eps}{\varepsilon}
\newcommand{\dd}{\,\mathrm{d}}
\begin{document}
\section*{UnsolvedMath ID: 700758}
\textbf{Source problem code:} catalogue entry
\subsection{Definitions and mathematical statement}
Let \(p\) be a prime. A pro-\(p\) group is a topological group isomorphic to an inverse limit of finite \(p\)-groups. Let \(F\) be a free pro-\(p\) group, with its prime and rank to be specified as discussed below.

Under the usual interpretation in the first-order language of groups \(\mathcal L_{\mathrm{grp}}=\{\cdot,{}^{-1},1\}\), define
\[
\operatorname{Th}_{\mathrm{grp}}(F)=\{\varphi:\varphi\text{ is an }\mathcal L_{\mathrm{grp}}\text{-sentence and }F\models\varphi\}.
\]
Here variables range over elements of the underlying group; the topology is not separately named in this language. This theory is decidable if there is an algorithm which, given any \(\mathcal L_{\mathrm{grp}}\)-sentence \(\varphi\), halts and determines whether \(F\models\varphi\).

\textbf{Question (source formulation).} Is the elementary theory of a free pro-\(p\) group decidable?

\textbf{Unresolved qualifications.} The supplied record does not specify the rank of the free pro-\(p\) group, whether the rank is finite or infinite, or whether trivial and procyclic cases are excluded. It also does not state whether \(p\) and the rank are fixed parameters, whether the question ranges over all such parameters, or whether a uniform decision algorithm is intended. The language above is the usual pure-group-language interpretation, which the source does not explicitly specify. These qualifications must be resolved before selecting a fully quantified version of the question.
\subsection{Short English statement}
The question asks whether an algorithm can decide every first-order sentence about a free pro-\(p\) group. The supplied formulation leaves the prime, rank, language, and any uniformity requirement unspecified. It names Yu. L. Ershov in the statement, while the proposer metadata names A. G. Myasnikov and V. N. Remeslennikov.
\subsection{Sources}
\begin{itemize}
\item Local UnsolvedMath registry, permanent ID \(700758\), source code \texttt{catalogue entry}, title \emph{Kourovka Notebook Problem 11.66}. The supplied record labels the problem open; this label is not an independently verified literature assessment.
\item \emph{The Kourovka Notebook (Unsolved Problems in Group Theory)}, original problem 11.66, 11th issue (1990); the supplied background identifies the 21st issue edition (2026), arXiv:1401.0300v48. The collection metadata names E. I. Khukhro and V. D. Mazurov as editors. This reference is supplied attribution and has not been independently consulted.
\item Attribution: the statement begins with Yu. L. Ershov's name in parentheses. The background and proposer metadata name A. G. Myasnikov and V. N. Remeslennikov, with proposed year 1990. The record does not explain the relationship between these attributions.
\end{itemize}
\end{document}
